\subsection{原函数的渐近展开}

设 E 是 \(+\infty\) 的一个邻域上的一个比较尺度，由 \(\neq 0\) 的、在 \(+\infty\) 的一个邻域上具有常号的实函数组成；设 f 是定义在一个区间 \([a, +\infty[\) 上的规则向量函数，取值于一个完备赋范空间 E 中，容许一个渐近展开

\[
\mathbf {f} = \sum_ {\lambda \leqslant \alpha} \mathbf {a} _ {\lambda} g _ {\lambda} + \mathbf {r} _ {\alpha}
\]

（精确到 \(g_{\alpha}\) ，相对于 E）。进一步假设每个原函数 \(\int_{a}^{x} g(t) dt\) （函数 \(g \in E\) 的）相对于 E 容许一个渐近展开。在这些情形下我们将看到，可以得到 \(\mathbf{F}(x) = \int_{a}^{x} \mathbf{f}(t) dt\) 相对于 E 的一个渐近展开。我们区分两种情形：

\(1^{\circ}\int_{a}^{+\infty}g_{\alpha}(t)dt\) 为无穷；则有 \(\int_{a}^{x}\mathbf{r}_{\alpha}(t)dt\ll \int_{a}^{x}g_{\alpha}(t)dt\) (V,p.230, prop. 6)；由假设可以得到 \(\sum_{\lambda \leqslant \alpha}\mathbf{a}_{\lambda}\int_{a}^{x}g_{\lambda}(t)dt\) 的一个渐近展开

（精确到某个 \(g_{\rho}\) ）(V, p.~222)；若 \(cg_{\sigma}\) 是 \(\int_{a}^{x} g_{\alpha}(t) dt\) 的主部，则由此将有 \(\int_{a}^{x} \mathbf{f}(t) dt\) 的一个精确到 \(g_{\min(\rho, \sigma)}\) 的渐近展开，其所有项的范数不定递增。

\(2^{\circ}\int_{a}^{+\infty}g_{\alpha}(t)dt\) 收敛；设 \(\beta\) 是指标 \(\lambda \leqslant \alpha\) 中最小的使 \(\mathbf{a}_{\lambda}\neq 0\) 且使 \(\int_{a}^{+\infty}g_{\lambda}(t)dt\) 收敛者；积分

\[
\mathbf {C} = \int_ {a} ^ {+ \infty} \left(\mathbf {f} (t) - \sum_ {\lambda <   \beta} \mathbf {a} _ {\lambda} g _ {\lambda} (t)\right) d t
\]

这时收敛，且可以写

\[
\mathbf {F} (x) = \sum_ {\lambda <   \beta} \mathbf {a} _ {\lambda} \int_ {a} ^ {x} g _ {\lambda} (t) d t + \mathbf {C} - \sum_ {\beta \leqslant \lambda \leqslant \alpha} \mathbf {a} _ {\lambda} \int_ {x} ^ {+ \infty} g _ {\lambda} (t) d t - \int_ {x} ^ {+ \infty} \mathbf {r} _ {\alpha} (t) d t.
\]

则 \(\int_{x}^{+\infty}\mathbf{r}_{\alpha}(t)dt\ll \int_{x}^{+\infty}g_{\alpha}(t)dt\) ；若 \(cg_{\sigma}\) 是 \(\int_{x}^{+\infty}g_{\alpha}(t)dt\) 的主部，且若有

\[
\sum_ {\lambda <   \beta} \mathbf {a} _ {\lambda} \int_ {a} ^ {x} g _ {\lambda} (t) d t + \mathbf {C} - \sum_ {\beta \leqslant \lambda \leqslant \alpha} \mathbf {a} _ {\lambda} \int_ {x} ^ {+ \infty} g _ {\lambda} (t) d t
\]

的一个渐近展开（精确到 \(g_{\rho}\) ），则由此将有 \(\mathbf{F}\) 的一个精确到 \(g_{\min(\rho, \sigma)}\) 的渐近展开。

于是问题归结为求 E 中函数的原函数相对于 E 的渐近展开。我们已经看到，在某些关于 E 的假设下，V, p.~233 的命题 8 如何给出这种原函数的主部。此外，命题 8 的证明给出 V, p.~233 的公式 (1) （相应地， (2)）两边之差的表达式，其形式为函数 \(\frac{1}{|\mu+1|}\left(xf'(x)+f(x)\right)-f(x)\) （相应地， \(f(x)g'(x)\) ，其中 \(g=f/f'\) ）的一个原函数；在构成这个新原函数的主部（作为 (1) （相应地， (2)）右边的一个渐近展开）时，就得到所求展开的右边（参见 V, p.~247-255）。

例. 1) 设 \(f(x) = 1 / \log x\) （ \(x > 1\) ）；我们已经看到 \(\int_{a}^{x} dt / \log t \sim x / \log x\) ，且差 \(\int_{a}^{x} \frac{dt}{\log t} - \frac{x}{\log x}\) 是 \(1 / (\log x)^2\) 的一个原函数；对这个函数可再次应用命题 8，从而得到 \(\int_{a}^{x} dt / (\log t)^2 \sim x / (\log x)^2\) 。由递归于是得到展开

\[
\begin{array}{l} \int_ {a} ^ {x} \frac {d t}{\log t} = \frac {x}{\log x} + \frac {x}{(\log x) ^ {2}} \\ \qquad + \frac {2 x}{(\log x) ^ {3}} + \dots + (n - 1)! \frac {x}{(\log x) ^ {n}} + o \left(\frac {x}{(\log x) ^ {n}}\right). \end{array}
\]

注意，不论 n 为何，这个展开的所有项都随 x 趋于 \(+\infty\) 。

\begin{enumerate}
\def\labelenumi{\arabic{enumi})}
\setcounter{enumi}{1}
\tightlist
\item
  设 \(f(x) = \frac{e^x}{x^2 + 1}\) ；可以写 \(f(x) = \frac{e^x}{x^2} - \frac{e^x}{x^4} + o_1\left(\frac{e^x}{x^4}\right)\) 。命题 8 给出展开
\end{enumerate}

\[
\int_ {a} ^ {x} \frac {e ^ {t}}{t ^ {2}} d t = \frac {e ^ {x}}{x ^ {2}} + \frac {2 e ^ {x}}{x ^ {3}} + \frac {6 e ^ {x}}{x ^ {4}} + o _ {2} \left(\frac {e ^ {x}}{x ^ {4}}\right), \quad \int_ {a} ^ {x} \frac {e ^ {t}}{t ^ {4}} d t = \frac {e ^ {x}}{x ^ {4}} + o _ {3} \left(\frac {e ^ {x}}{x ^ {4}}\right)
\]

由此，把两式相加，

\[
\int_ {a} ^ {x} \frac {e ^ {t}}{t ^ {2} + 1} d t = \frac {e ^ {x}}{x ^ {2}} + 2 \frac {e ^ {x}}{x ^ {3}} + 5 \frac {e ^ {x}}{x ^ {4}} + o _ {4} \left(\frac {e ^ {x}}{x ^ {4}}\right).
\]

